\subsection{Estimation results}

The financial cycle indicator is derived from a dynamic factor model with a single factor and four lags. The model is estimated separately for each country, yielding country-specific measures of the financial cycle. The single-factor specification is motivated by the objective of capturing the dominant co-movement among the variables included in the model. In practice, the estimated factor is closely related to the first principal component of the data, which explains the largest share of their joint variation.

\paragraph{Adjustment to the data.}

Figure \ref{var_exp} reports the proportion of total variance explained by the first factor for each country. The dashed horizontal line represents the cross-country average, showing that the first factor accounts for approximately 47\% of the total variance on average. In several countries, this share exceeds 50\%, indicating that a single-factor representation provides a reasonably good summary of the common dynamics in the data.

\begin{figure}[htbp]
    \centering
    \includegraphics[width=1\linewidth]{TABLES_FIGURES/var_exp.png}
    \caption{\% of total variance explained by the first factor}
    \caption*{\footnotesize \textbf{Notes:} The figure reports the share of total variance explained by the estimated factor of the dynamic factor model estimated separately for each country. The dashed horizontal line indicates the cross-country average. All variables are expressed in 2-year growth rates (2-year differences for bond spreads), filtered to remove low-frequency components, and standardized prior to estimation.}
    \label{var_exp}
\end{figure}

Table \ref{table_r2} reports the $R^2$ values obtained by regressing each variable on the estimated common component, for all countries in the sample. These values provide a simple measure of how closely each variable co-moves with the financial cycle indicator. The results show that the estimated factor explains a large share of the variation in credit to non-financial corporations, the debt service ratio, and the credit-to-GDP ratio. This finding suggests that the common factor is primarily driven by variables related to credit and indebtedness, which are central to the build-up of cyclical systemic risk.

By contrast, the explanatory power of the factor is generally lower for share prices and sovereign bond spreads, indicating that these variables co-move less systematically with the rest of the dataset. One interpretation is that these series are more strongly influenced by country-specific shocks, global factors, or abrupt changes in market sentiment, which are not fully captured by a single common factor. This pattern is broadly consistent with \citet{drehmann2012characterising}, who emphasize the central role of credit variables in characterizing the financial cycle. At the same time, our results suggest that house prices play a more heterogeneous role across countries and, on average, contribute less to the common factor than credit-based indicators.

Some very low $R^2$ values also stand out in Table \ref{table_r2}, particularly for credit to households and house prices in Belgium, and for house prices in Portugal and Sweden. These cases likely reflect the fact that, after the transformations and filtering applied to the data, the corresponding series display relatively weak co-movement with the other variables entering the model. In such cases, a single common factor naturally attributes little explanatory power to those series. This does not imply that these variables are unimportant for financial stability in those countries; rather, it indicates that their cyclical dynamics are less aligned with the dominant common component identified by the model.

\begin{table}[htbp]
  \centering
  \caption{$R^2$ of the first factor across variables}
  \begin{adjustbox}{width=1\textwidth}
    \begin{tabularx}{\textwidth}{lXXXXXXX}
    \toprule
    \textbf{Country} & \textbf{Credit NFC} & \textbf{Credit HH} & \textbf{House prices} & \textbf{Share prices} & \textbf{DSR} & \textbf{Credit-to-GDP} & \textbf{Bond spread} \\
    \midrule
    BE    & 0.89  & 0.03  & 0.00  & 0.26  & 0.90  & 0.87  & 0.08 \\
    DE    & 0.66  & 0.21  & 0.14  & 0.14  & 0.80  & 0.80  & 0.05 \\
    ES    & 0.92  & 0.84  & 0.36  & 0.04  & 0.82  & 0.79  & 0.09 \\
    FI    & 0.81  & 0.59  & 0.12  & 0.14  & 0.87  & 0.77  & 0.56 \\
    FR    & 0.84  & 0.50  & 0.19  & 0.24  & 0.70  & 0.68  & 0.11 \\
    GB    & 0.74  & 0.55  & 0.38  & 0.13  & 0.73  & 0.76  & 0.14 \\
    IT    & 0.87  & 0.78  & 0.66  & 0.24  & 0.74  & 0.74  & 0.26 \\
    NL    & 0.41  & 0.18  & 0.22  & 0.01  & 0.81  & 0.66  & 0.02 \\
    PT    & 0.88  & 0.31  & 0.01  & 0.17  & 0.60  & 0.69  & 0.09 \\
    SE    & 0.82  & 0.51  & 0.01  & 0.37  & 0.78  & 0.81  & 0.33 \\
    \midrule
    \textbf{Average} & \textbf{0.78} & \textbf{0.45} & \textbf{0.21} & \textbf{0.17} & \textbf{0.77} & \textbf{0.76} & \textbf{0.17} \\
    \bottomrule
    \end{tabularx}
  \end{adjustbox}
  \label{table_r2}
  \parbox{\textwidth}{\raggedright \footnotesize \textbf{Notes:} The table reports the $R^2$ from regressing each variable on the estimated common factor of the dynamic factor model, computed separately for each country. All variables are expressed in 2-year growth rates (2-year differences for bond spreads), filtered to remove low-frequency components, and standardized prior to estimation. Higher values indicate stronger co-movement with the common factor.}
\end{table}

\paragraph{Evolution of the financial cycle across countries.}

Figure \ref{factor_median} reports the median financial cycle indicator across the countries in the sample (results for individual countries are provided in Appendix \ref{AppendixA}). Estimates are obtained using three alternative methods: principal components (PCA), quasi-maximum likelihood (QML), and the two-step estimator (TSTEP). The figure also displays the interquartile range (25th–75th percentiles), capturing the cross-country dispersion of the financial cycle across estimators. The similarity of the estimated cycles across methods indicates that the main features of the financial cycle are robust to the choice of estimator.

\begin{figure}[htbp]
    \centering
    \includegraphics[width=1\linewidth]{TABLES_FIGURES/factor_medians.png}
    \caption{Median of the financial cycle across countries}
    \caption*{\footnotesize \textbf{Notes:} The figure shows the median financial cycle indicator across countries, together with the interquartile range (25th–75th percentiles). Results are reported for three estimation methods: principal components (PCA), quasi-maximum likelihood (QML), and two-step estimator (TSTEP). The shaded area represents cross-country dispersion across all estimators.}
    \label{factor_median}
\end{figure}

The results show that the financial cycle is characterised by medium-term fluctuations, with phases lasting on the order of 8 to 12 years. This distinguishes the financial cycle from the business cycle, which typically operates at higher frequencies. Over the sample period, the financial cycle exhibits three prominent peaks: in the early 2000s, in the run-up to the GFC in 2007–2008, and following the COVID-19 shock. These episodes coincide with periods of strong credit expansion, rising leverage, and increasing systemic risk.

Following the GFC, the financial cycle enters a prolonged subdued phase, characterised by lower amplitude and reduced variability. This pattern is consistent with the post-crisis regulatory environment, which includes tighter capital requirements, constraints on leverage, and the broader use of macroprudential policy tools aimed at containing the build-up of systemic risk.

The post-COVID period is marked by a sharp but short-lived expansion of the financial cycle, reflecting the unprecedented scale of fiscal and monetary support measures. This phase is associated with rapid credit growth, abundant liquidity, and rising asset prices. However, from 2022 onwards, the financial cycle contracts as inflationary pressures intensify and monetary policy tightens globally. The resulting increase in interest rates leads to a deterioration in financing conditions and a slowdown in credit dynamics.

Overall, these results highlight the close interaction between macroeconomic conditions, policy interventions, and the financial cycle. In particular, they illustrate how expansionary policies can amplify financial dynamics, while monetary tightening and regulatory constraints can contribute to their reversal.

\subsection{Decomposition}

The factor estimates can be further decomposed into contributions from the underlying variables, providing a clearer understanding of the main drivers of the financial cycle. This decomposition involves performing a linear regression of the estimated factor (obtained using the PCA estimator) on the variables included in the model. The regression coefficients are converted to absolute values and normalized to sum to one, with each variable's contribution calculated as the product of its normalized coefficient and observed value.

Figure \ref{dec_graph} illustrates the decomposition of the median of the financial cycle indicator across the set of countries in the sample (results for individual countries are available in Appendix \ref{AppendixB}). {\color{blue}[\textit{There is something in these decomposition graphs that I do not understand. The bars should in principle sum to the black line, am I right? But in most periods it looks like they do not. Are you omitting the constant? If so we should say it otherwise it looks wrong. I think you could also redistribute the constant contribution to the other variables, in equal parts, for representation purposes. In both cases, the graph needs a note explaining what we did. I would also specify somewhere that the decomposition is performed in-sample (am I right?)}]} The financial cycle shows a pronounced build-up in the early 2000s, primarily driven by strong growth in household credit and non-financial corporations credit, alongside rising DSR and credit-to-GDP ratios. This expansion is followed by a sharp contraction during the mid-2000s, as credit growth decelerates and the financial cycle indicator approaches its long-term average.

A significant resurgence occurs leading up to the 2007–2008 GFC, largely due to an increase in the growth of credit to households and increasing house prices. The onset of the GFC triggers a sharp and prolonged downturn in the financial cycle, a decline further exacerbated by the European sovereign debt crisis in 2010. During this period, the contribution of most variables - including credit indicators and house prices - turn negative, reflecting widespread deleveraging and financial distress across European economies.

In the second half of the 2010s, the financial cycle stabilizes at relatively subdued levels, as credit growth remains weak despite gradual improvements. Notably, house prices emerge as a more significant driver of the financial cycle during this phase, reflecting growing overvaluation concerns in certain European housing markets. A renewed recovery is observed following the COVID-19 shock, driven by credit expansion and increased liquidity from unprecedented fiscal and monetary support.

Around 2022, the financial cycle enters a renewed downturn as global inflation accelerates and central banks implement rapid monetary tightening. The decline in the credit-to-GDP and house prices weighs significantly on the financial cycle.

Overall, the decomposition underscores that while credit growth and leverage measures remain the primary drivers of financial cycle fluctuations, asset prices and interest rate spreads have played an significant role during specific periods. These insights can enhance the monitoring of financial cycles, help identify emerging trends and risks, and ultimately inform macroprudential policy.

\begin{figure}[htbp]
    \centering
    \includegraphics[width=1\linewidth]{TABLES_FIGURES/dec_graph.png}
    \caption{Decomposition of the median of the financial cycle across countries}
    \caption*{\footnotesize \textbf{Notes:} The figure shows the decomposition of the median financial cycle indicator across countries into contributions from the underlying variables. Contributions are based on a linear projection of the estimated factor on the standardized variables and are normalized to sum to the total variation of the indicator.}
    \label{dec_graph}
\end{figure}

\subsection{Identification of neutral risk environment}

\paragraph{Our approach.} 

A neutral risk environment corresponds to a state in which cyclical systemic risk is neither elevated nor subdued. It is typically associated with moderate credit growth, stable asset price dynamics, and financial conditions that do not signal the build-up or materialisation of imbalances. In such periods, vulnerabilities evolve in line with underlying economic fundamentals, without clear signs of overheating or stress.

To operationalize this concept, we adopt a quantitative approach based on the financial cycle indicator. Specifically, we define a neutral regime as periods in which the indicator remains close to its long-term average. Drawing on the asymptotic distribution theory of the principal components estimator for the common factor, we construct confidence bands around the estimated financial cycle. These bands provide a statistical benchmark to identify significant deviations from typical conditions. Periods in which the indicator lies within the confidence bands are classified as neutral, while deviations outside these bands signal either elevated or subdued cyclical systemic risk.

\paragraph{Empirical evidence.} 

Figure \ref{ci_medians} illustrates this approach using the median financial cycle indicator across countries (country-specific results are reported in Appendix \ref{AppendixC}). The shaded area represents the 80\% confidence bands, while the dashed lines denote the upper and lower confidence bounds. This framework allows for a clear classification of financial cycle phases. Elevated regimes correspond to periods in which the indicator exceeds the upper bound, signalling heightened systemic risk. Subdued regimes are identified when the indicator falls below the lower bound, typically reflecting the materialisation of risk and subsequent deleveraging. Neutral regimes correspond to periods in which the indicator remains within the confidence bands.

For the median across countries, the results identify elevated financial cycle phases in the early 2000s and in the run-up to the GFC, coinciding with strong credit expansion and rising asset prices. In contrast, the period following the European sovereign debt crisis is characterised by a subdued regime, reflecting weak credit growth and subdued financial activity. From around 2016 onwards, the financial cycle gradually converges towards a neutral regime. As of 2024, the most recent observations suggest that the financial cycle has moved back into a subdued phase.

\begin{figure}[htbp]
    \centering
    \includegraphics[width=1\linewidth]{TABLES_FIGURES/ci_medians.png}
    \caption{Financial cycle and 80\% confidence intervals}
    \caption*{\footnotesize \textbf{Notes:} The figure displays the median financial cycle indicator across countries together with 80\% confidence bands based on the asymptotic distribution of the estimated factor. The shaded area represents the confidence bands, while the dashed lines denote the upper and lower bounds. Values above (below) the bands indicate statistically significant deviations associated with elevated (subdued) cyclical systemic risk.}
    \label{ci_medians}
\end{figure}

\paragraph{Cross-country analysis of financial cycle regimes.} 

Extending this framework across countries, Figure \ref{regimes_graph} classifies financial cycle conditions into subdued, neutral, and elevated regimes over time. Subdued regimes are prevalent in the mid-to-late 1990s and during the post-GFC period, reflecting deleveraging and weak credit dynamics. Elevated regimes are concentrated in the early 2000s and in the run-up to the GFC, consistent with the build-up of systemic risk. Following the COVID-19 shock, many countries temporarily transition towards neutral conditions, supported by expansionary fiscal and monetary policies. Since 2022, however, several countries have returned to subdued regimes as tighter monetary policy and higher interest rates weigh on credit growth and financial activity.

\begin{figure}[htbp]
    \centering
    \includegraphics[width=1\linewidth]{TABLES_FIGURES/regimes_graph.png}
    \caption{Financial cycle regimes across countries}
    \caption*{\footnotesize \textbf{Notes:} The figure classifies financial cycle conditions across countries into subdued, neutral, and elevated regimes based on the position of the financial cycle indicator relative to its estimated confidence bands. Each bar represents the share of countries in each regime at a given point in time.}
    \label{regimes_graph}
\end{figure}

\paragraph{Policy implications.} 

The calibration of the countercyclical capital buffer (CCyB) is subject to considerable uncertainty, as the build-up of cyclical systemic risk is difficult to measure in real time. In practice, this raises the risk of delayed or insufficient policy action, as well as the possibility of underestimating vulnerabilities during periods of apparent stability.

In response to these challenges, the conce  pt of a positive neutral CCyB has gained prominence. As of November 2025, 17 macroprudential authorities in the European Economic Area (EEA) have implemented a positive neutral CCyB framework. A positive neutral CCyB refers to a capital buffer maintained when risks are assessed as neither elevated nor subdued, ensuring that banks hold releasable capital in normal times. Evidence from the COVID-19 episode shows that the release of such buffers supports lending and mitigates procyclical deleveraging, alleviating concerns about banks’ reluctance to use capital buffers \citep{avezum2024assessment, berrospide2021used}.

Our approach enables the identification of periods in which cyclical systemic risks are neither building up nor materialising. By assessing deviations in the financial cycle, constructed as deviations of underlying indicators from their long-term averages, it provides a benchmark for evaluating when risks lie within a neutral range. In this context, the framework offers a transparent and operational tool to inform the calibration the neutral level of the CCyB. 

\subsection{Relation with Financial crisis}

{\color{blue} [\textit{Consider renaming this subsection as ``Identification of financial crises", to make it coherent with the previous one}]}

As documented by \citet{borio2014financial}, one defining feature of financial cycles is that they tend to be closely associated with systemic banking crisis. This link suggests that a good measure of the financial cycle can be used as a leading indicator to signal the build-up of risk that may culminate in financial crisis. In this section, we investigate the early warning properties of our proposed measure of the financial cycle. 

Our analysis is conducted using the systemic banking crisis dataset constructed by \cite{lo2017new} for European countries. A crisis, as defined by this dataset, refers to a period during which the financial system served as a source or amplifier of economic shocks, financial intermediaries faced significant distress or collapse, and substantial crisis management policies were implemented. 

\paragraph{Set-up of the early warning analysis.} We estimate country-specific logit regressions to assess the build-up of systemic risk ahead of banking crises. The dependent variable is a binary vulnerability indicator that takes the value one during the five to twelve quarters preceding a systemic banking crisis and zero otherwise. This definition emphasises the objective of detecting early signs of financial imbalances, rather than predicting the exact timing of a crisis. This distinction is important, as macroprudential policy must be preemptive to be effective. Moreover, there are implementation lags that constrain policy responsiveness. For example, the CCyB includes a one-year phase-in period, allowing banks time to adjust to new capital requirements. As a result, signals issued less than a year before a crisis are likely too late for an effective policy response. The focus is therefore on identifying vulnerabilities early enough to enable timely intervention.

As independent variable, we consider the financial cycle measure obtained in the context of a dynamic factor model estimated with different estimators. Our analysis is purely univariate as we do not consider additional variables our lags of the indicator. To benchmark our results, we use the Basel gap, calculated as the cyclical component of the total credit-to-GDP ratio derived from a recursive HP filter with a smoothing parameter of 400,000. The Basel gap is chosen as the benchmark indicator because it is widely considered as a leading univariate signal of systemic financial crises. It serves serves as the recommended reference indicator for activating the countercyclical capital buffer under the guidance of the Basel Committee on Banking Supervision.

We use two measures of accuracy in the early warning exercise: the area under the operating receiver curve (AUROC); and the usefulness index proposed by \cite{alessi2011quasi}. The AUROC is a global measure of accuracy and plot relates the ratio of false positives (noise ratio) to the ratio of true positives (signal ratio) for every possible value for threshold that can be applied to the predicted probabilities coming from the logit regression. \citet{hanley1982meaning} show that the AUROC can be simply interpreted as the probability of correctly classifying zeros and ones. 

The relative usefulness {\color{blue}\st{measure} index [\textit{we are talking about the usefullness \textbf{index} here, right? If so, let's call it in a uniform way throughout the text, it is clearer.}]} represents a measure of adjustment that takes into account the preference of policymakers regarding type one (not signaling a financial crisis, denoted by $T1$) and type two (``false alarms", denoted by $T2$) errors. The indicator is defined as follows:

\begin{equation}
    U = \frac{min[\theta; 1- \theta] - L(\theta)}{min[\theta; 1- \theta]},
\end{equation}

where $L(\theta)$ denotes the policymakers loss function which is defined as

\begin{equation}
   L(\theta) = \theta T1 + (1- \theta) T2.
\end{equation} 

The policymakers preference are controlled by $\theta$. A value for $\theta$ lower than 0.5 indicates that the policymaker is is less averse towards missing a signal of a crisis than receiving a false alarm. Similarly to \citet{lang2019anticipating} we consider balance preferences for type one and type two errors. Because the policy maker can always realise a loss of $min[\theta; 1- \theta]$ by disregarding the indicator, the usefulness {\color{blue}measure [\textit{same comment as above}]} compares the loss of following the logit model results against the loss of disregarding it. 

Finally, we consider the full sample used in the previous analysis, with two exceptions. First, Finland is excluded from the early warning exercise, as it did not experience a systemic banking crisis during the available sample period. Second, we exclude the crisis periods themselves from the sample used to estimate the logit models. This follows the recommendation of \citet{bussiere2006towards}, who show that early warning model performance can be distorted by the disorderly and volatile adjustment of explanatory variables following the onset of a crisis. To avoid this so-called “post-crisis bias,” we focus only on the pre-crisis and tranquil periods in the estimation.

\paragraph{In-sample exercise.} The results for all countries, presented in Table \ref{ew_in_sample}, show that the financial cycle indicator constructed using a dynamic factor model has good early warning properties. On average, it consistently outperforms the credit-to-GDP gap in terms of the AUROC, while maintaining a comparable usefulness index. Among the different estimators, the PCA-based indicator performs slightly better than both the QML and two-step estimators.

\begin{table}[htbp]
  \centering
  \caption{In-sample exercise}
  \small % Reduce font size
  \begin{adjustbox}{width=\textwidth}
    \begin{tabularx}{\textwidth}{l p{0.1cm} X X X X p{0.25cm} X X X X}
    \toprule
          &       & \multicolumn{4}{c}{\textbf{AUROC}} &       & \multicolumn{4}{c}{\textbf{Usefulness index}} \\
    \cmidrule{3-6}\cmidrule{8-11}
    \textbf{Country} &       & \textbf{PCA} & \textbf{QML} & \textbf{TSTEP} & \textbf{BG} &       & \textbf{PCA} & \textbf{QML} & \textbf{TSTEP} & \textbf{BG} \\
    \midrule
    BE    &       & 0.54  & 0.56  & 0.55  & 0.58  &       & 0.20  & 0.20  & 0.21  & 0.29 \\
    DE    &       & 0.65  & 0.50  & 0.43  & 0.46  &       & 0.17  & 0.14  & 0.06  & 0.18 \\
    ES    &       & 0.94  & 0.85  & 0.91  & 0.51  &       & 0.45  & 0.40  & 0.43  & 0.23 \\
    FR    &       & 0.45  & 0.63  & 0.57  & 0.62  &       & 0.10  & 0.23  & 0.17  & 0.31 \\
    GB    &       & 0.61  & 0.52  & 0.59  & 0.64  &       & 0.27  & 0.23  & 0.28  & 0.32 \\
    IT    &       & 0.90  & 0.98  & 0.92  & 0.53  &       & 0.37  & 0.47  & 0.37  & 0.26 \\
    NL    &       & 0.68  & 0.49  & 0.61  & 0.57  &       & 0.25  & 0.15  & 0.20  & 0.29 \\
    PT    &       & 0.52  & 0.60  & 0.56  & 0.51  &       & 0.23  & 0.25  & 0.26  & 0.26 \\
    SE    &       & 0.65  & 0.53  & 0.67  & 0.65  &       & 0.17  & 0.13  & 0.18  & 0.32 \\
    \midrule
    \textbf{Avg} &       & \textbf{0.66} & \textbf{0.63} & \textbf{0.65} & \textbf{0.56} &       & \textbf{0.25} & \textbf{0.24} & \textbf{0.24} & \textbf{0.27} \\
    \bottomrule
    \end{tabularx}%
  \end{adjustbox}
  \label{ew_in_sample}%
    
    \parbox{\textwidth}{\raggedright \footnotesize \textbf{Notes:} PCA, QML, and TSTEP refer to estimates of the financial cycle indicator obtained using the principal components, quasi-maximum likelihood, and two-step estimators, respectively. BG denotes the credit-to-GDP gap (Basel gap). AUROC is the area under the receiver operating characteristic curve. The usefulness index is computed following \citet{alessi2011quasi} under balanced preferences between type I and type II errors.}
  
\end{table}%

\paragraph{Simulated out-of-sample exercise.} We also conduct an out-of-sample exercise focusing on the GFC to assess the early warning properties of the financial cycle indicator. Evaluating such properties for this period is particularly challenging due to the very low frequency of systemic banking crisis. Despite this constraint, the credibility of an early warning model depends on its ability to detect, in real time, the buildup of financial imbalances preceding a crisis. 

To address this issue, we follow an approach similar to \cite{alessi2018identifying} and check the model's predictions as of mid-2006, using only data available up to the second quarter of that year. Crucially, this exercise excludes any knowledge of the subsequent GFC {\color{blue}\st{around 2008}}, ensuring a realistic assessment of the model’s practical utility. To expand the scope of countries included in the out-of-sample exercise, we incorporate ``residual events" as defined in \cite{detken2014operationalising}. These events represent periods {\color{blue}\st{where} in which}, absent policy interventions or external factors that dampened the credit cycle, a systemic financial crisis {\color{blue}\st{tied} related} to the domestic financial cycle would likely have occurred. {\color{blue}[\textit{Please see if my correction makes sense; it took me a while to understand this passage, so please consider rephrasing to make it easier to read.}]}

The exercise is conducted in pseudo-real time, meaning that all transformations are based solely on the information available up to each point in time. However, we acknowledge two simplifying assumptions: first, we rely on final revised data rather than vintages, and second, we assume no publication lag, meaning observations for a given period are assumed to be available immediately. The procedure involves estimating the dynamic factor model using data up to 2006 Q1. We then perform a logit regression of the vulnerability indicator on the financial cycle indicator and use the resulting model to predict the probability of a financial crisis occurring within one year. This process is repeated for each quarter through 2009 Q1.

\begin{figure}[htbp]
    \centering
    \includegraphics[width=1\linewidth]{TABLES_FIGURES/oos_graph.png}
    \caption{Simulated out-of-sample predicted probability of a financial crisis in 1 year}
    \label{oos_graph}
    \caption*{\footnotesize \textbf{Notes:} The figure reports the median predicted probability of a systemic banking crisis within one year across countries, together with the interquartile range. Predictions are obtained from country-specific logit models estimated in pseudo real-time using only information available up to each point in time.}
\end{figure}

Figure \ref{oos_graph} illustrates the cross-country median of the predicted crisis probabilities alongside the 25th–75th interquartile range. The results show that the median estimated probability of a crisis within one year rises significantly from 10\% in 2006 Q1 to around 25\% in 2007. The results show that the model would successfully produce a signal of future distress in the financial system. These findings suggest that even a simple univariate early warning model based on the financial cycle indicator could provide macroprudential authorities with valuable guidance and reasonable advance notice. This constitutes encouraging evidence for the indicator’s potential usefulness in policy settings.

{\color{blue}[\textit{One could argue that 25\% probability is still pretty low. Let's caveat this result, acknowledging the objective difficulties of a model that has never faced a crisis like the GFC to anticipate it. Moreover, we could express it in relative terms to stress the increase in probability, like "The results show that the estimated median probability of a crisis within one year more than doubles between 2006 and 2007, increasing significantly from 10\% in the first quarter of 2006 to around 25\% in 2007."}]}
